\section{Related Work}
\label{sec:related_work}
\subsection{BEV Perception}
BEV provides a unified space for integrating information from multiple cameras, making it one of the primary representations used in autonomous driving.
A crucial step in BEV perception lies in transforming features from the perspective camera space into the BEV space.
In this section, we categorize existing transformation approaches into three main groups: \textit{learning-based unprojection}, \textit{geometry-based unprojection} and \textit{projection from learned 3D volumes}.

\noindent\textbf{Geometry-based unprojection.}
Geometry-based approaches explicitly lift 2D image coordinates into 3D space.
Lift-Splat-Shoot (LSS) \cite{philion2020lift} predicts per-pixel depth and uses camera intrinsics and extrinsics to map image features to 3D points through a geometric unprojection process.
However, its performance heavily depends on the quality of the predicted depth.
Since the training process relies solely on downstream task supervision, the depth estimation can be unstable or inaccurate.
Subsequent works \cite{li2023bevstereo, li2023bevdepth, xie2022m, yang2023parametric} extend the LSS paradigm by introducing additional depth supervision to improve the accuracy and robustness of depth prediction.
BEVDepth \cite{li2023bevdepth} leverages LiDAR-derived depth maps for supervision, while BEVStereo \cite{li2023bevstereo} employs stereo matching techniques to enhance depth estimation quality.

\noindent\textbf{Learning-based unprojection.}
Learning-based methods \cite{liu2022petr, liu2023petrv2, pan2023baeformer, zhou2022cross, bartoccioni2023lara} implicitly learn the mapping from 2D to 3D coordinates in a data-driven manner. 
These approaches employ attention-based architectures or multi-layer perceptrons (MLPs) to directly associate image features with 3D spatial locations.
Instead of explicitly predicting depth, they encode 3D positional information into learnable queries and use cross-attention to aggregate multi-view image features into BEV space.
However, these methods remain fully implicit, lacking geometric reasoning and interpretability.

\noindent\textbf{Projection from learned 3D volumes.}
Instead of directly lifting image features or learning implicit mappings, another category of methods constructs a predefined 3D volume and learns dense feature representations within it.
BEVFormer \cite{li2024bevformer} and \cite{harley2023simple} learn a 3D feature grid that captures spatial structure and semantics across the scene.
To obtain BEV features, these methods project or collapse the learned 3D feature volume onto the BEV plane, generating top-down representations.
PointBEV \cite{chambon2024pointbev} further introduce a coarse-to-fine training strategy that focuses on regions of interest to reduce memory consumption

However, these methods are trained solely with downstream task losses without explicit 3D reconstruction, leading to a lack of interpretability.
In contrast, we propose Splat2BEV, which introduces explicit 3D reconstruction during both training and inference, and projects the learned representation into BEV space for downstream tasks.

\subsection{Gaussian Splatting in Autonomous Driving}
Recently, a growing number of studies have incorporated Gaussian Splatting into the autonomous driving tasks.

\noindent\textbf{Navigation.}
Several recent works \cite{keetha2024splatam, matsuki2024gaussian, yan2024gs, huang2024photo, peng2024rtg} integrate 3D Gaussian Splatting into robot navigation and SLAM systems as the mapping component.
They extend traditional 3D Gaussian Splatting to support online reconstruction, enabling continuous scene updates during exploration.
These approaches demonstrate that Gaussian primitives can provide a compact and differentiable representation for mapping in autonomous driving scenarios.

\noindent\textbf{Simulation.}
Another line of research \cite{yan2024street, zhou2024drivinggaussian, zhou2024hugs, hess2025splatad} leverages 3D Gaussian Splatting to reconstruct realistic environments for autonomous driving simulation.
They employ Gaussian Splatting to build photorealistic and geometrically consistent driving scenes from real-world data.
By transforming large-scale outdoor environments into editable Gaussian representations, these approaches enable efficient rendering, controllable scene manipulation, and sensor simulation for training and evaluating perception and planning models in autonomous driving.

\noindent\textbf{BEV Perception.}
GaussianLSS \cite{lu2025toward} and GaussianBEV \cite{chabot2025gaussianbev} introduces 3D Gaussian Splatting as an intermediate representation, enabling a more efficient aggregation of BEV feature.
However, these two Gaussian Splatting–based methods treat the Gaussians merely as feature carriers rather than performing explicit 3D scene reconstruction.
Their learning process still relies solely on downstream task supervision, limiting geometric interpretability.

